En el análisis de datos funcionales, las series se convierten en funciones
mediante una expansión en bases, un paso cuyo efecto rara vez se examina en
aplicaciones económicas. Esta tesis compara
las bases B-spline, de Fourier y de ondículas (Haar, Daubechies~2, Symlet~4) para
representar el crecimiento anual del producto interno bruto de 20 países
latinoamericanos entre 1960 y 2024, extendiendo un estudio previo para incluir
la contracción por COVID-19. Las bases se
comparan mediante el error de reconstrucción, el número de coeficientes no nulos
con error equivalente y el análisis de componentes principales funcionales. Los
datos extendidos reproducen la varianza del estudio original (47.2\% frente a
47.1\% en tres componentes). Con error
equivalente, el B-spline adaptativo utiliza menos coeficientes que las ondículas
(19.5 nodos frente a aproximadamente 38 de 128 coeficientes), y ambas familias
distribuyen el error de forma distinta: los B-splines siguen los episodios
prolongados pero suavizan los choques aislados, mientras que las ondículas
umbralizadas conservan algunos choques aislados pero atenúan los episodios
amplios. La base de Fourier es la menos adecuada, pues las series no muestran
una periodicidad clara y sus aproximaciones oscilan en periodos estables. La
energía por niveles de ondículas identifica las escalas de variación de cada
país. El número de componentes necesario para explicar el
80\% de la varianza varía entre tres y ocho según la base y el suavizado, lo que
muestra que la elección de la base es una decisión de modelización que
condiciona los resultados del análisis.\\

\noindent \textbf{Palabras Clave}:\\
Análisis de datos funcionales, expansión en bases, ondículas, B-splines,
componentes principales funcionales, crecimiento económico, América Latina.
